% Part II: the Databases tab. Sources: docs/user-guide.md ("The Databases tab"); docs/databases.md;
% medorama/web/js/views/databases.js; medorama/src/databases.rs; gate.rs (column_mask).
\chapter{The Databases tab}
\label{chap:databases-tab}

The \ui{Databases} tab shows the same data as the patient pages, organized by database. It lists
every database in use, with what it holds, and shows each one as a table of its rows as they
were imported, every row linked to its patient and its file to its viewer. Databases are the
only way data enter Medorama: a file (a recording, an ECG, an image) always comes with the row
that names its patient. \Cref{chap:importing} describes how databases are imported, and
\cref{chap:connecting} how they are connected.

\screenshot{databases.png}{The Databases tab: every database in use, one of them unfolded into
its rows.}{fig:databases-tab-list}

\section{The list of databases}
\label{sec:databases-tab-list}

The list shows every database in use, newest first. Three kinds of databases appear in it:
\begin{itemize}
\item \textbf{imported tables} (CSV, Parquet or spreadsheet files, copied into the catalog), one
  database per table, or per set of files imported together;
\item \textbf{folders of recordings}, marked \ui{folder};
\item \textbf{connected databases} (an OMOP instance, a warehouse, read in place), marked
  \ui{connected} with their kind and engine (for example \ui{connected $\cdot$ OMOP $\cdot$
  PostgreSQL}).
\end{itemize}

Each line gives the following:
\begin{description}
\item[\ui{Database}] The database's name: the file's name, or the name given when it was
  imported or connected. Cohort criteria refer to it by this name (\cref{chap:cohorts}).
\item[\ui{Holds}] The categories of data that its rows filled (\cref{chap:patients}).
\item[\ui{Rows}] Its number of rows (for a connected database, \ui{live}).
\item[\ui{Patients}] The number of patients it holds data about.
\item[\ui{Source}] The source it was imported under, when that differs from its own name.
\item[\ui{Imported}] When it was imported, and by whom.
\end{description}
When there are some, the line also gives the rows \ui{held back} for review and the items
\ui{waiting} for a database that names their patient (\cref{chap:importing}).

\paragraph{Unfolding a database.} A click on a database's line unfolds its table in place (the
arrow at the start of the line turns), and a second click folds it again. Several tables can be
open at once, and the open ones are remembered in this browser. Two kinds of lines do not
unfold:
\begin{itemize}
\item A connected database, whose line says ``its data are read from it when a patient's page
  opens or a cohort is searched''. Its persons, stays and entries appear on the patients' pages
  (\cref{chap:patients}).
\item A table imported before Medorama kept rows, whose line says ``Imported before rows were
  kept: import it again to see it here''.
\end{itemize}

\paragraph{A page of its own.} \ui{Full page}, in an unfolded table, shows that database alone,
at an address of its own (\code{\#/databases/12}). The page shows its categories, its numbers
of rows and patients, its source, and when and by whom it was imported, with a search line
across the page (``Search this database's rows''). The address keeps the search, the filters,
the order and the page, so that the view can be reloaded or shared.

\section{A database's table}
\label{sec:databases-tab-table}

The rows are shown \emph{as they were imported}: the file's own columns, under their own names,
in their own order, a hundred rows to a page (\ui{$\langle$ Previous}, \ui{Next $\rangle$}).
Two columns come before the file's columns:
\begin{description}
\item[\ui{Line}] The row's line in the file.
\item[\ui{Patient}] The patient the row reached: the patient ID (in Privacy Mode, the
  pseudonym), as a link to the patient's page at the entry the row made. A dash means that no
  database names that patient yet.
\end{description}
The last column holds the row's link to its file (\cref{sec:databases-tab-files}) and, when the
row is not an ordinary one, its status: \ui{patient, encounter} (a row that describes who the
patient is, or an encounter), \ui{duplicate}, \ui{held back}, \ui{not read} or
\ui{nothing to read}. Rows held back or not read are dimmed.

The tooltip of a column's header tells what the column was mapped to when it was imported
(``$\rightarrow$ \code{mrn}'', ``$\rightarrow$ \code{value}'' and so on), or ``not mapped''. A
long cell is cut after 80 characters, and its tooltip shows it whole.

\paragraph{Searching, filtering, sorting.}
\begin{itemize}
\item \textbf{Search.} The box above the table (``Search these rows''), or the search line of
  the table's own page, keeps the rows in which every word is found in some cell, ignoring
  case.
\item \textbf{Filters.} The box under each column's name (\ui{filter}) keeps the rows whose
  cell in that column contains the text, ignoring case. Filters combine with each other and with
  the search.
\item \textbf{Order.} A click on a column's name sorts the rows by that column, and a second
  click sorts them in the other direction; an arrow shows which. Numbers are sorted as numbers,
  then text, with empty cells last, whichever the direction and whichever database (SQLite or
  PostgreSQL) holds the catalog. Without a sort, the rows are in the file's order. The search
  and the filters take \code{\%} and \code{\_} as written.
\end{itemize}

\section{Rows, patients and files}
\label{sec:databases-tab-files}

Each row links to \emph{its patient}, and its file to its viewer. In the file's column, the cell
shows the file's name, followed by one of these links:
\begin{description}
\item[\ui{VitalStream}] For a recording (ICU monitoring, a Holter, an EEG, a sleep study),
  which opens in the VitalStream tab at its start, for a minute. When the catalog links to a
  viewer running elsewhere, the recording opens in that viewer's browser tab.
\item[\ui{ECG}] For a 12-lead ECG, which opens in the ECG tab.
\item[\ui{PACS}] For a DICOM study, a series, an echo loop, a picture, a video or a PDF, which
  opens in the PACS tab.
\item[\ui{Open}] For any other file, which the browser opens in a new tab.
\end{description}
A viewer opened from a row says whose file it shows, as it does from the patient's page
(\cref{chap:viewers-in-medorama}). The item's category and the file's kind decide which viewer
opens which file (\cref{chap:ecg-imaging}).

\section{Identifying columns}
\label{sec:databases-tab-phi}

\begin{phibox}
A database's columns are of three kinds, according to how its import mapped them:
\begin{itemize}
\item \textbf{Mapped columns that do not identify} (a patient ID, an encounter ID, a time, a
  test's name, a value, a unit and so on), and \textbf{measurement columns} (a column named like
  a measurement with its unit, such as \code{QT (ms)} or \code{LVEF (\%)}, or a known ECG
  measurement; \cref{chap:cohorts}). These are shown to everyone.
\item \textbf{Identifying columns}: those mapped to an MRN, its authority, a last or first name,
  a date of birth or an encounter (account) number\phimark.
\item \textbf{Unrecognized columns}: columns that are not mapped and are not measurements.
\end{itemize}
Out of Privacy Mode, users with the PHI permission see, search, filter and sort every column.
In Privacy Mode, every column is still listed under its imported name, so that the table keeps
the shape of its source, but its cells are sent as follows:
\begin{itemize}
\item the cells of identifying columns become placeholders (\code{[MRN]},
  \code{[Patient Name]}, \code{[Date of Birth]}, \code{[CSN]}), patient IDs become
  \code{[Patient ID]}, file names \code{[File Name]} and the authors of notes
  \code{[Staff Name]};
\item the cells of unrecognized columns become \code{[Withheld]}, while measurement columns are
  shown;
\item the text of notes is shown, unless an administrator withholds \ui{The text of notes and
  reports}, in which case it becomes \code{[Withheld]};
\item the other cells pass the privacy gate: encounter IDs become pseudonyms, and dates move
  back.
\end{itemize}
The search, the filters and the sort use the cells as you are shown them: pseudonyms in place
of identifiers, each patient's dates moved back, notes without the patient's names. The search
skips placeholders and \code{[Withheld]} cells (and the text of notes, when it is withheld); a
filter on a placeholder column compares with the placeholder, and a sort by one of those columns
is ignored, so that they cannot reveal what is withheld. Such searches and filters are written
to the audit log (\code{privacy\_search}). An administrator chooses what Privacy Mode withholds
(\cref{chap:users-privacy}); the rules are explained in \cref{chap:protecting}.
\end{phibox}

\section{Exporting and deleting (administrators)}
\label{sec:databases-tab-admin}

Administrators have three more buttons:
\begin{description}
\item[\ui{Import databases\dots}] Above the list, this button opens the \ui{Import} tab
  (\cref{chap:importing}).
\item[\ui{Export CSV}] In a table, this button exports the database as CSV, with its columns
  under their imported names and its rows as imported. In Privacy Mode, the button reads
  \ui{Export CSV (Privacy Mode)}, and the export has the same shape but passes the gate:
  identifying cells become placeholders, patients are named by their pseudonyms, and dates move
  back. A folder of recordings, or a table imported before rows were kept, has no rows to
  export.
\item[\ui{Delete\dots}] In a table, this button removes the database from the catalog. Medorama
  first says what would go (so many recordings, patients, encounters, bed stays and items, and
  what stays because another database also holds it). It then asks whether to delete the
  catalog's indexes and results of the removed recordings as well (in the results folder only).
  Files are never touched.
\end{description}

\paragraph{Changes show everywhere at once.} When a database is imported or deleted, here or in
anyone else's browser, the change shows within seconds on every open page: the list, the
patients' pages, the cohorts' searches, and the viewers, which close when their file is gone
(\cref{chap:page}). \Cref{chap:importing} describes how to remove data and sources from the
command line.
